% s-161-p161-m002; original PDF 161, 162
\Aphorism{s-161-p161-m002}{247}
Ich zähle das Glück des Lebens nach Punkten; Linien, Flächen und Körper kommen für mich nur insoweit in Betracht, als sie sich in Punkten schneiden können. Ich danke dem Leben nicht für seine langen, breiten und planmässigen Veranstaltungen, für unangefochtene Jugend, für Dauer von Neigungen und Freundschaften, für ruhiges Parallellaufen mit Menschen, Büchern, Zeitbegebnissen; das Alles ist werthvoll, aber es ist zu vernünftig, um zu begeistern. Aber ich bin auf das Leben gut zu sprechen, weil es neben dem Kosmos auch ein Chaos hat, und in diesem Chaos die wunderlichsten Zufälle, Augenblicke, Zusammenklänge — ich blicke in einen Hexenkessel, worin immer wieder einmal eine teufelsmässige Delicatesse gebraut wird. Es sind nur Punkte, Begegnungen, auf Messers Schneide tanzende Wunder, glatte schlüpfrige unwahrscheinliche Dinge, die Einem leicht wie ein Orangenkern aus den Fingern gleiten — jedenfalls aber waren sie einmal zu fassen, und diese \Foreign{points} eben zählen im Spiel! Ich danke dem Leben, wofür? für jene blitzenden Augenblicke zwischen zwei Welttheilen und zwei Jahrhunderten, für Höhen-, Umkehr- und Wendepunkte, für Pässe und Bergrücken, für Grenzen, Brücken und Übergänge, für aufleuchtende und verschwindende Regenbögen, für alles Einmal-und-nicht-wieder, — ich danke ihm, mythologisch gesprochen, dafür, dass man auf Secunden und nicht mehr als Secunden die Accorde der \Sperr{prästabilirten Harmonie} zu hören bekommt! „Wenig ist die Art des besten Glücks“. Suchen wir nach dem Unverlierbaren in unserem Gedächtniss, nach den Erinnerungspuren genossener Lust, so finden wir irgend etwas Kleines,Flüchtiges, Nichtssägliches, die Mücke im Bernstein: vielleicht ein Glas \Foreign{Asti spumante,} ein Wölkchen am Himmel, Blüthenduft, Blick, Gebärde, den Tonfall einer Phrase oder einen Tact Gartenmusik, irgend ein Mehr-als-Nichts und Weniger-als-Etwas, das der Seele Lust und Wohlgeschmack am Leben giebt und ihre Begehrlichkeit mehr reizt als sättigt: ein \Foreign{pianissimo} von Glück, auf das hin es sich lohnt zu leben und noch nicht lohnt zu sterben \Dots{3} Giebt es etwas Haftenderes als diese Wellenspiele und unfassbaren Übergänge, diese zitternden Gleitungen und Schwebungen der Seele? Mir ist alles Dauernde, wie Person, Landschaft, Kunstwerk, nur werthvoll als Möglichkeit und Reservoir solcher Augenblicke. Ich frage mich, wieviel ich etwa an einer Person höchstens gewinnen oder verlieren kann, ja ob ich der Empfindung des Verlustes überhaupt fähig bin: bleiben nicht die Perlen, wenn die Perlenschnur zerreisst? Jemand, der mir auch nur einen jener leuchtenden Augenblicke geschenkt hat, ist dadurch für mich unzerstörbar, unsterblich geworden. Ich kenne keine andere Unsterblichkeit als die des Gewesenseins — meine Schatzkammer birgt nichts Köstlicheres als diese Tropfen, die zu Edelsteinen erstarrt sind.
\par

